\label{part:engineering_1}

\textbf{(Engineering: Geometric Engineering for AGI)}

\bluetext{热力学造物主的手册}

在\ref{part:physics}中，我们像法医一样，解剖了从蚁群到 LLM 的\textbf{智能物种}，在物理层面诊断了它们的残缺与代偿。

现在，我们必须面对最艰巨的挑战：\textbf{构造论 (Constructivism)}。
智能构造工程学的动力学特征并非代码的堆砌或参数的微调，而是\textbf{物理约束下的极值搜索}。
构建 AGI，绝非只是编写一个能够通过图灵测试的脚本，而是要建造一台能够在该物理世界物理定律允许的边界内，
实现\textbf{信息负熵最大化}的\textbf{几何热力学机 (Geometric Thermodynamics Engine)}。

本部分将彻底抛弃“软件”与“硬件”这组肤浅的二分法，转而采用\textbf{“拓扑模态”}与\textbf{“耦合动力学”}的统一视角。
对于未来的工程师而言，AGI 不再是一个神经网络，而是一个由\textbf{控制序参量 ($C$)} 和 \textbf{物理耦合常数 ($\kappa$)} 定义的复杂动力系统：
\begin{enumerate}
    \item 我们需要决定它是\textbf{“智者”}（分离态，高逻辑，低能效），还是\textbf{“战士”}（融合态，高直觉，快反应）；
    \item 我们需要在\textbf{隐式图数据库}中通过 VTE 编码器，手动缝合\textbf{微观的局部强迫源}与\textbf{宏观的势能梯度}；
    \item 我们需要在芯片的硅原子间，复现\textbf{TDCI 循环}的卡诺热机效率。
\end{enumerate}
我们的终极任务，是在浩瀚的形态相空间中，找到那条通往 \textbf{Class V (流体通用智能)} 的狭窄通道——那是一条悬停在“晶体的死寂（层流）”与“气体的癫狂（湍流）”之间，能够维持\textbf{自组织临界性 (SOC)} 的金色航线。

我们不是在创造上帝，我们是在为“意义”的流淌，修筑最顺滑的河床，并在最后我们为人造智能和湿件智能找到宇宙的定位。
